This paper showed, in theory and in experiments, that a single fixed LED and a single photodiode reoriented to a few known orientations determine the three-dimensional position of the receiver. The measurement model separates a direction, given by the shape of the RSS over the orientations, from a distance, given by their common amplitude through the calibrated gain and pattern of the LED. The Fisher information makes the conditions explicit: the receiver must be tilted with respect to its rotation axis, the slope of its angular response sets the attainable accuracy, and orientation errors act as noise that averaging cannot remove. On these grounds the testbed used a $15^\circ$ tilt and the \del{measured}\add{published} response model of the photodiode\add{, confirmed by the shape test}.

On 125 positions over five heights, the calibrated system reached a mean 3D error of 3.66~cm with ten orientations and 3.74~cm with five; on 131 blind random positions estimated live with five orientations, it reached 4.10~cm. The same accuracy is obtained with 10~ms of signal per orientation instead of 100~ms, and the computation of a position then takes about 7~ms. Each element of the model mattered, from the receiver response (77~cm with a Lambertian receiver) to the angular zero of the rotation stage and the position of the LED, and each can be calibrated with little effort: nine flat-PD points over the working area calibrate the LED as well as a dense map, provided they cover the emission angles in use. The comparison with the position error bound showed that the remaining error is due to model bias, not to noise; the residual of the fit, available without ground truth, detected faults of the orientation stage; and the same data revealed the position and tilt of every LED in the room. Better models of the LED pattern and of the receiver orientation, faster reorientation, and the combination of several LEDs are the next steps.
